\chapter{Bounds and complementarity}\label{chap:constraints}

A model with only equalities cannot represent a limit that binds in some
scenarios and not in others: an import quota, a capacity limit, a ceiling on
borrowing. Such limits are complementarity conditions, set with
\manref{cge_constraint} and solved with the rest of the model, so that the
solver finds which limits bind.

\section{Rents}\label{sec:rents}

Three rent variables give quotas and capacity limits their price effect. Each
is zero in the base and fixed by every closure, unless a constraint frees it
(section~\ref{sec:cons-def}):

\begin{cgetable}{>{$}l<{$}p{0.82\linewidth}}
\toprule
\text{Variable} & Description\\
\midrule\endhead
\bottomrule\endlastfoot
\V{RM}_c & premium on imports of $c \in CM$, per unit of the world price in local currency\\
\V{RE}_c & premium on exports of $c \in CE$, per unit of the world price\\
\V{RA}_a & rent on the output of activity $a$, per unit of the activity price\\
\end{cgetable}

They enter the price equations \eqref{eq:P1}, \eqref{eq:P2} and \eqref{eq:P7}.
A premium is a wedge between the world price and the domestic price, like a
tariff, but its revenue is a rent. The rent on each element,
\begin{equation*}
R^M_c = \V{RM}_c\,\V{pwm}_c\,\V{EXR}\,\V{QM}_c, \qquad
R^E_c = \V{RE}_c\,\V{pwe}_c\,\V{EXR}\,\V{QE}_c, \qquad
R^A_a = \V{RA}_a\,\V{PA}_a\,\V{QA}_a,
\end{equation*}
is paid to one recipient $\varrho$: the government (the default, as when quota
licences are auctioned), a household or enterprise, a factor, or the rest of
the world (as with a voluntary export restraint, whose rents foreign exporters
keep). Rents are added to the incomes in \eqref{eq:I1} (factors),
\eqref{eq:I3} (households and enterprises) and \eqref{eq:I8} (government), and
rents paid abroad to the payments abroad in \eqref{eq:S3}, at $R/\V{EXR}$. A
premium fixed at a positive value (\cmd{cge_shock RM}) acts as a tariff whose
revenue goes to the recipient.

\textbf{GDP.} The income side adds the rents received at home and every
capacity rent, which is value added produced at home wherever it is paid. The
expenditure side values imports and exports net of the rents paid abroad.

\section{Constraints}\label{sec:cons-def}

A constraint $k$ pairs a limited variable $Y_k$ (one element of a variable,
for example $\V{QM}_{com\_man}$) and its bound $b_k$ with an adjusting variable
$X_k$, which the closure fixes at $\bar{x}_k$ and which may move in one
direction only:
\begin{equation}
g_k = \sigma_k\,\frac{b_k - Y_k}{s^Y_k} \ge 0, \qquad
h_k = \delta_k\,\frac{X_k - \bar{x}_k}{s^X_k} \ge 0, \qquad
g_k\,h_k = 0
\tag{B1}\label{eq:B1}
\end{equation}
with $\sigma_k = 1$ for an upper bound and $-1$ for a lower bound, and
$\delta_k = 1$ if $X_k$ may rise and $-1$ if it may fall. The scales $s$ are
the base magnitudes of the variables (one where the base value is zero).
Either $Y_k$ is at its bound and $X_k$ has moved from $\bar{x}_k$ to keep it
there, or $X_k$ is at $\bar{x}_k$ and $Y_k$ is within its bound. Each
constraint adds one equation and frees one variable, so the system stays
square.

\begin{cgetable}{p{0.22\linewidth}p{0.16\linewidth}p{0.16\linewidth}p{0.38\linewidth}}
\toprule
Constraint & Limited variable & Adjusting variable & Effect when it binds\\
\midrule\endhead
\bottomrule\endlastfoot
import quota & $\V{QM}_c \le b$ & $\V{RM}_c$, up & the premium rises until imports equal the quota\\
export quota & $\V{QE}_c \le b$ & $\V{RE}_c$, up & the premium rises until exports equal the quota\\
capacity & $\V{QA}_a \le b$ & $\V{RA}_a$, up & the activity price rises above cost; the rent goes to the owners of the capacity\\
borrowing limit under a peg & $\V{FSAV} \le b$ & $\V{EXR}$, up & the peg holds until foreign savings reach the limit; the currency then depreciates\\
deficit ceiling & $\V{GSAV} \ge b$ & $\V{GADJ}$, down & government consumption is cut once the deficit reaches the ceiling\\
deficit ceiling, taxes & $\V{GSAV} \ge b$ & $\V{TINSADJ}$, up & direct tax rates rise instead\\
nominal wage floor & $\V{W}_f \ge b$ & $\V{QFS}_f$, down & employment falls once the wage reaches the floor\\
\end{cgetable}

$Y_k$ must be free and $X_k$ fixed under the closure. With $X_k$ a rent
variable, the constraint and the rent's recipient are given together.

\section{Solution}\label{sec:cons-solution}

From version 0.14.0 every condition is written in the Fischer--Burmeister
form
\begin{equation*}
\phi(g_k, h_k) = g_k + h_k - \sqrt{g_k^2 + h_k^2} = 0,
\end{equation*}
which holds exactly where $g_k \ge 0$, $h_k \ge 0$ and $g_k h_k = 0$
\citep{fischer1992}, and the conditions and the rest of the model are solved
at once by the semismooth Newton method of \citet{deluca1996}. $\phi$ is not
differentiable where $g_k = h_k = 0$, but it is semismooth: each step solves
the system of an element of its generalised Jacobian, computed from the exact
derivatives, $\partial\phi/\partial g = 1 - g/\sqrt{g^2 + h^2}$ and
$\partial\phi/\partial h = 1 - h/\sqrt{g^2 + h^2}$, with $1 - 1/\sqrt{2}$ for
both at $g = h = 0$. The step $d$ is taken along a line search on the merit
function $\psi(z) = \tfrac12 \lVert F(z)\rVert^2$, with $F$ the scaled
residuals, when it is a direction of sufficient descent,
$\nabla\psi' d \le -10^{-8}\lVert d\rVert^{2.1}$; otherwise, or when its line
search fails, the step is the steepest descent of the merit function,
$-\nabla\psi = -J'F$. The line search is non-monotone \citep{grippo1986}: a
step $t\,d$ is accepted when $\psi(z + t d) \le \max_{j \le 10} \psi_{k-j} +
10^{-4}\,t\,\nabla\psi' d$, the largest merit of the last ten iterations,
which the method needs on hard problems such as that of Kojima and Shindo. Near
a solution where every constraint is either binding or not, with its other
condition strictly satisfied, the method converges as Newton's method does.

Once the system is solved, each constraint takes the regime of the solution:
binding, $g_k = 0$, if $h_k > g_k$, and otherwise not binding, $h_k = 0$; the
model is solved again in these regimes, which polishes the solution, unless a
regime's other condition then fails. Should the semismooth method fail, the
active-set method of version 0.4.4 follows: the smoothed form
\begin{equation*}
g_k + h_k - \sqrt{g_k^2 + h_k^2 + 2\mu} = 0, \qquad \mu = 10^{-8},
\end{equation*}
finds the regimes, and the model is solved in them, a constraint whose other
condition fails changing regime, for at most ten rounds;
\cmd{cge_options, mcp(activeset)} (\manref{cge_options}) makes it the method.

The solution therefore satisfies the conditions exactly, to the solver's
tolerance. The test suites check the method on problems of MCPLIB
\citep{dirkse1995} with known solutions: that of \citet{josephy1979} and that
of \citet{kojima1986}, with a degenerate solution, from several starts; on
an economy in the form of MPSGE whose technologies are used or not, against
the solution of its linear programme; and on many constraints of the
standard model at once, against the active-set method and PATH. Three
implementations check one another, each with its own method: Stata, with the
methods above; the Python reference, with an active-set
method that starts from the solution without constraints; and GAMS, as a mixed
complementarity problem solved with PATH (\secref{sec:gams-types}). The
minimum wage of chapter~\ref{chap:labour} keeps its smoothed condition and its
calibration.

\section{Reporting}

\manref{cge_solve} reports each constraint's regime, and the regimes are stored
with the simulation. The rent variables appear in the results like other
variables. Constraints are kept when the shocks are reset, in dynamic runs and
sensitivity analysis, and when the model is saved.

\section{User inequalities}

From version 0.9.0, \manref{cge_equation} writes inequalities among any
expressions of the user language, each paired with a variable fixed under the
closure (\secref{sec:userineq}). Each element is a constraint of this chapter
on a slack variable, the difference between the two sides, bounded at zero
and adjusted by the paired variable, and the bounds of user variables are
constraints on the variables themselves. They are listed and solved with the
constraints of \manref{cge_constraint}.

\section{Marginals}\label{sec:cons-marginals}

From version 0.9.8, the solution listing of \manref{cge_listing},
\opt{sol}{ution}, gives each constraint as GAMS lists a complementarity pair:
its level is the room within its bound, $g_k$ in the units of $Y_k$, and its
marginal is the move of its adjusting variable, $h_k$ in the units of $X_k$,
such as the premium of a quota or the fall in \cmd{GADJ} under a deficit
ceiling. One of the two is zero. \manref{cge_gams} writes the scenario with the
adjusting variables at their fixed levels, so that PATH finds the regimes
itself, and GAMS reports the same room and marginal for each pair.

\section{Bounds that do not bind}

The bounds of \manref{cge_bound} (from version 0.9.8) are not constraints. As
the bounds of a GAMS CNS model, they keep the solver within a domain and are
not meant to bind: the system stays square, with no variable to adjust. A
solution beyond such a bound cannot be reached; a limit that binds is a
constraint of this chapter.

